\scanpage{50}
(pour pouvoir passer à la $\varprojlim$ dans la catégorie des schémas), pour $V$ un schéma $\varprojlim$ de schémas séparés réguliers de t.f. sur $\mathbb{Q}$ (morphismes de transition \inserted{\unsure{imm. ouvertes}} affines sans plus), Alors les morphismes \uline{dominants} de schémas $V\longrightarrow U$ doivent correspondre aux hom. ext $E_V\longrightarrow E_U$, compatibles avec les projections dans $E_{\mathbb{Q}}\simeq\Gamma_{\mathbb{Q}}$, et \unsure{telles} que l'image soit ouverte. P.ex. on pourrait prendre \inserted{pour} $U$, \deleted{$\operatorname{Spec}(\unsure{\mathcal{O}}_{X,x})$}, $V=$ les spectres d'anneaux locaux\inserted{ réguliers}, \unsure{ess.|ets.} de type fini sur $\mathbb{Q}$.

\figure{\begin{tikzpicture}\draw (0,0)--(.55,.03);\end{tikzpicture}}

Cette conjecture fondamentale (éventuellement revue et corrigée en cours de route !) étant admise, la question qui se pose ensuite est de déterminer les topos (pro) galoisiens sur $B_{\mathbb{Q}}$ qui proviennent de modèles élémentaires anabéliens -- ou encore, les $\Pi_U=\pi_1(\bar U)$ de tels modèles étant connus, de déterminer quelles sont
\margin{49}
